When large language models help users write messages, they make implicit choices along pragmatic dimensions---formality, directness, politeness, emotional tone, and specificity---that shape the communicative impact of the resulting text. Yet it remains unclear which of these dimensions the model actively ``considers'' during generation versus which it can reason about only when explicitly prompted. We introduce \emph{perplexity sensitivity analysis}, a behavioral probing method that detects pragmatic variables under consideration by measuring whether an LLM assigns significantly different likelihoods to response variants that differ along a single pragmatic axis. Across 30 message-writing scenarios spanning workplace, personal, and service domains, we find that two instruction-tuned models (\qwenseven and \qwenthree) exhibit strong, statistically significant directional preferences for four of five axes: polite over blunt ($d{=}1.81$), vague over specific ($d{=}1.02$), indirect over direct ($d{=}0.94$), and formal over casual ($d{=}0.68$). Emotional tone is the sole axis with no detectable preference. Comparing these behavioral results against model introspection reveals a striking dissociation: directness is strongly defaulted on but rarely self-reported, while emotional tone is frequently mentioned but not behaviorally detected. These results replicate perfectly across model sizes (Spearman $\rho{=}1.0$), suggesting that perplexity sensitivity analysis is a reliable, low-cost method for auditing pragmatic defaults in LLMs.
